% ============================================================================
% Detection depth of spectral defects in shifted zeta strings
% Draft v1 (2026-08-29), produced from MANUSCRIPT_defect_depth_theorem_proof_skeleton.md
% Sources of record: LAB_ihara_round{1..5}.md, ROUND6_closeout.md, and the two
% companion papers (shifted screw-function margin paper; shifted zeta-string paper).
% NOTE: bibliography entries are drafted from the lab record and MUST be verified
% against the actual sources before submission (marked %% VERIFY).
% ============================================================================
\documentclass[11pt]{amsart}
\usepackage{amsmath,amssymb,amsthm}
\usepackage[margin=1.15in]{geometry}
\usepackage{booktabs}
\usepackage{xcolor}
\definecolor{linkblue}{rgb}{0.1,0.15,0.5}
\usepackage[colorlinks=true,linkcolor=linkblue,citecolor=linkblue,urlcolor=linkblue]{hyperref}
\usepackage{microtype}

\newtheorem{theorem}{Theorem}[section]
\newtheorem{proposition}[theorem]{Proposition}
\newtheorem{lemma}[theorem]{Lemma}
\newtheorem{corollary}[theorem]{Corollary}
\theoremstyle{definition}
\newtheorem{definition}[theorem]{Definition}
\newtheorem{assumption}[theorem]{Assumption}
\theoremstyle{remark}
\newtheorem{remark}[theorem]{Remark}
\newtheorem{openproblem}[theorem]{Open problem}

\newcommand{\R}{\mathbb{R}}
\newcommand{\C}{\mathbb{C}}
\newcommand{\dd}{\,d}
\newcommand{\muh}{\widehat{\mu}}
\newcommand{\eps}{\varepsilon}
\DeclareMathOperator{\sgn}{sgn}
\DeclareMathOperator{\supp}{supp}
\DeclareMathOperator{\Rea}{Re}
\DeclareMathOperator{\Ima}{Im}

\begin{document}

\title[Detection depth of spectral defects in shifted zeta strings]
{Detection depth of spectral defects\\ in shifted zeta strings}

\author{Edward B. Baker III}
\date{Draft of August 29, 2026}

\begin{abstract}
We study how a finite-rank spectral defect is detected by a positive
Stieltjes--Weyl system. After a resolvent compactification of an unbounded
Stieltjes measure at a scale $s_0>0$, the compactified moments are an
invertible binomial transform of the Taylor coefficients of the Weyl function
at the negative resolvent point $-s_0$. Finite-rank defects perturb the
associated Hankel determinants through exact Christoffel--Darboux kernel
formulas, and the first loss of Stieltjes positivity is located exactly at the
first anomalous coefficient of the associated Stieltjes continued fraction.
For a fixed exterior defect over a regular compactly supported background the
detection order satisfies $n_\eps/\log(1/\eps)\to (2g)^{-1}$, where $g$ is the
Green function of the complement of the admissible support. We apply this
framework to the positive shifted zeta-string reference family associated with
Suzuki's shifted screw functions: a hypothetical zero quartet beyond the
shifted boundary enters as a conjugate rank-two defect with effective
displacement $\eta=\delta-\omega$, the same parameter that controls
exponential growth in the time-domain criterion, and the signed
absolutely continuous correction it produces is dominated by the reference
density and cannot affect the leading detection depth. We derive the
resulting inverse-spectral sensitivity law and compare it with the
time-domain detection scale
$\eta t_*=2\log\tau+\log\log\tau+\log(S_\omega/\eta)+o(1)$.
The results quantify the very different conditioning of two exact analytic
representations of the same hypothetical defect; they provide neither a new
zero-free region nor any progress toward the Riemann hypothesis.
\end{abstract}

\maketitle
\tableofcontents

% ============================================================================
\section{Introduction}\label{sec:intro}
% ============================================================================

\subsection{The question}
A positive measure $\sigma$ on $[0,\infty)$ with
$\int (1+\lambda)^{-1}\dd\sigma(\lambda)<\infty$ corresponds, through its
Stieltjes--Weyl function
\[
q(z)=\int_0^\infty\frac{\dd\sigma(\lambda)}{\lambda-z},
\qquad z\in\C\setminus[0,\infty),
\]
to a positive inverse-spectral object: a Kre\u{\i}n string, equivalently a
positive Stieltjes continued fraction \cite{Stieltjes,KacKrein,Wall,GK}.
If the spectral data are corrupted by a small \emph{defect} located outside
the admissible support --- a point mass on the negative axis, or a conjugate
pair of non-real point masses --- then positivity of the inverse problem must
eventually fail. The question this paper quantifies is:
\begin{quote}
\emph{How much finite-resolution spectral information is required before a
small off-support defect becomes detectable as a failure of positivity?}
\end{quote}
Three classical finite-resolution coordinates are available: truncated moment
sequences, Taylor jets of the Weyl function at a fixed point of the resolvent
set, and initial segments of the Stieltjes continued fraction (equivalently,
initial pieces of the string). A first contribution of this paper is that,
after a natural one-parameter \emph{resolvent compactification} that makes all
moments of an unbounded Stieltjes measure finite, these three coordinates are
exactly equivalent, degree by degree
(Theorem~\ref{thm:compactification}). The central mechanism is then:
\[
\boxed{\;
\text{spectral defect}
\;\longrightarrow\;
\text{finite-rank Hankel perturbation}
\;\longrightarrow\;
\text{Christoffel threshold}
\;\longrightarrow\;
\text{first CF failure}.\;}
\]
Finite-rank perturbations of Hankel matrices are controlled exactly by
Christoffel--Darboux kernels (Theorems~\ref{thm:rankone}
and~\ref{thm:ranktwo}); the location of the first failed positivity condition
translates exactly, through the classical determinant representation of the
Stieltjes continued fraction, into the index of the first anomalous
continued-fraction coefficient (Theorem~\ref{thm:cfindex}); and for a fixed
exterior defect over a regular compactly supported background, exterior
polynomial growth turns the threshold into the asymptotic law
\begin{equation}\label{eq:introdepth}
\frac{n_\eps}{\log(1/\eps)}\longrightarrow\frac1{2g(c)}
\qquad(\eps\downarrow0),
\end{equation}
where $g$ is the Green function of the complement of the admissible support
with pole at infinity (Theorem~\ref{thm:depth}). We call $n_\eps$ the
\emph{detection depth} of the defect: it is simultaneously a moment order, a
Taylor-jet order, and (up to the exact parity rule of
Theorem~\ref{thm:cfindex}) the position of the first anomalous bead of the
associated string.

\subsection{The arithmetic application}
Suzuki's shifted screw functions \cite{Suzuki2023} attach to each shift
$\omega>0$ a criterion function $\Psi_\omega$ whose eventual nonnegativity is
equivalent to a zero-free half-plane
$\Rea\rho\le\tfrac12+\omega$ for the Riemann xi function, and, in the
companion papers \cite{MarginPaper,StringPaper}, a Weyl function
\begin{equation}\label{eq:qomega}
q_\omega(z)=\frac1{w}\,\frac{\xi'}{\xi}\Big(\tfrac12+\omega+w\Big),
\qquad w=\sqrt{-z},\ \Rea w>0 ,
\end{equation}
which, when the corresponding half-plane is zero-free --- in particular under
the Riemann hypothesis --- is a Stieltjes function whose spectral measure is
absolutely continuous with an explicit Poisson-broadened zero density. This
\emph{positive shifted-zeta reference system} is a structured arithmetic
background for the abstract theory: a hypothetical zero
$\rho=\tfrac12+\delta\pm i\tau$ with $\delta>\omega$ contributes to
\eqref{eq:qomega} a conjugate pair of principal-sheet poles at
\[
z_\pm=-(\eta\pm i\tau)^2,\qquad \eta:=\delta-\omega>0,
\]
each of Stieltjes weight $2$ (Proposition~\ref{prop:quartet}) --- a rank-two
defect in the sense of the abstract theory, whose distance from admissibility
is controlled by exactly the \emph{horizontal excess} $\eta=\delta-\omega$
that also drives the exponential instability
$e^{\eta t}$ of the time-domain criterion. The signed absolutely continuous
correction that accompanies the poles is dominated pointwise by the reference
density and cannot influence the leading detection depth
(Lemma~\ref{lem:cut}). The resulting sensitivity law is
Theorem~\ref{thm:zetadepth}; the time-domain counterpart, with its constant
term identified, is
\begin{equation}\label{eq:introtime}
\eta\,t_* \;=\; 2\log\tau+\log\log\tau+\log\frac{S_\omega}{\eta}+o(1),
\qquad
S_\omega=\frac{\xi'}{\xi}\Big(\tfrac12+\omega\Big),
\end{equation}
(Theorems~\ref{thm:envelope} and~\ref{thm:bracket}).

\subsection{What is, and is not, claimed}
Three caveats govern the entire paper. First, the zeta application is a
\emph{sensitivity statement about the positive reference system perturbed by
one hypothetical quartet}; if the Riemann hypothesis fails, the actual
response need not decompose as a positive background plus a single finite-rank
defect, and we never assume that it does (see the discussion before
Theorem~\ref{thm:zetadepth}). Second, the compactification is a natural
\emph{one-parameter} family indexed by the scale $s_0>0$; the detection order
is cutoff-independent for fixed $s_0$ but genuinely $s_0$-dependent, and no
claim of unique canonicity is made (Section~\ref{sec:s0}). Third, the
high-height law $n\asymp \tau^2\log\tau/(\eta\sqrt{s_0})$ is a moving-point
asymptotic: it follows from the fixed-point theorem only formally, and we
present it at the level actually established --- an explicit geometric
computation of the Green distance plus numerical confirmation in the regimes
accessible to computation (Section~\ref{sec:highheight}).
The results are statements about the \emph{conditioning of two exact
representations} of the same hypothetical object. They provide no new
zero-free region and no evidence that the inverse-spectral representation
brings the Riemann hypothesis closer.

\subsection{Relation to existing work}
Every individual algebraic ingredient below is classical: the Stieltjes
moment problem and continued fraction \cite{Stieltjes,Wall,Schmudgen}, the
Hankel determinant representation of the continued-fraction coefficients
\cite{Wall,GK}, finite-rank (Uvarov) modifications of orthogonality measures
\cite{Uvarov,Zhedanov}, Christoffel functions and Christoffel--Darboux
kernels \cite{Nevai,Totik}, regularity and $n$-th root asymptotics
\cite{StahlTotik}, and transformations between Stieltjes and Hausdorff moment
problems \cite{BergDurea}. The principle that Christoffel-function growth
detects points outside the support is also actively used in applied
anomaly detection \cite{LasserrePauwels,ChristoffelAnomaly2026,DataStreams}.
What we believe is new is the synthesis: the exact identification of the
first failed positivity condition with the first anomalous continued-fraction
coefficient, giving a precise inverse-spectral meaning to detection depth;
the use of the resolvent compactification as a finite-moment/Taylor-jet depth
coordinate for unbounded Weyl data, with its exact transform and length
invariance; and the application to the shifted zeta-string family, where the
defect parameter coincides with the horizontal excess $\eta=\delta-\omega$ of
the time-domain criterion. A final literature audit for these specific
combinations accompanies the submission.

% ============================================================================
\section{Stieltjes preliminaries and resolvent compactification}
\label{sec:compact}
% ============================================================================

Throughout, $\sigma$ denotes a positive Borel measure on $[0,\infty)$ with
\begin{equation}\label{eq:stieltjescond}
\int_0^\infty\frac{\dd\sigma(\lambda)}{1+\lambda}<\infty,
\end{equation}
and $q(z)=\int_0^\infty(\lambda-z)^{-1}\dd\sigma(\lambda)$ its
Stieltjes transform, holomorphic on $\C\setminus[0,\infty)$. We work with
this normalization (no additive affine term); the statements below extend to
the general Stieltjes normalization by trivial modifications.

\begin{definition}[Resolvent compactification at scale $s_0$]
\label{def:compact}
Fix $s_0>0$ and set
\[
x=\frac{\lambda}{s_0+\lambda}\in[0,1),
\qquad
\zeta=\frac{z}{s_0+z},
\qquad
z=\frac{s_0\zeta}{1-\zeta}.
\]
Let $\mu_{s_0}$ be the pushforward under $\lambda\mapsto x$ of the measure
$\dd\sigma(\lambda)/(s_0+\lambda)$, and let
\[
Q_{s_0}(\zeta)=\int_0^1\frac{\dd\mu_{s_0}(x)}{x-\zeta}.
\]
\end{definition}

\begin{theorem}[Resolvent compactification and Taylor-jet/moment equivalence]
\label{thm:compactification}
Let $\sigma$ satisfy \eqref{eq:stieltjescond} and let $s_0>0$. Then:
\begin{enumerate}
\item $\mu_{s_0}$ is a finite positive measure on $[0,1]$; in particular all
compactified moments $\muh_k=\int_0^1 x^k\dd\mu_{s_0}(x)$ are finite.
\item The transforms are related by
\begin{equation}\label{eq:Qtransform}
q(z)=(1-\zeta)\,Q_{s_0}(\zeta),
\qquad\text{equivalently}\qquad
Q_{s_0}(\zeta)=\frac{s_0+z}{s_0}\,q(z).
\end{equation}
\item Writing $q(z)=\sum_{j\ge0}c_j(z+s_0)^j$ near $z=-s_0$, one has
$c_j=\int_0^\infty(s_0+\lambda)^{-j-1}\dd\sigma(\lambda)$ and
\begin{equation}\label{eq:binomial}
\muh_k=\sum_{j=0}^k\binom{k}{j}(-s_0)^j c_j .
\end{equation}
In particular, for every $N$, the moments $(\muh_k)_{k\le N}$ and the Taylor
coefficients $(c_j)_{j\le N}$ determine each other through the invertible
triangular transform \eqref{eq:binomial}.
\end{enumerate}
\end{theorem}

\begin{proof}
(1) is \eqref{eq:stieltjescond} after the change of variables, since
$\dd\mu_{s_0}$ has total mass $\int (s_0+\lambda)^{-1}\dd\sigma\le
\max(1,s_0^{-1})\int(1+\lambda)^{-1}\dd\sigma<\infty$.

(2) A direct computation gives
\[
x-\zeta=\frac{\lambda}{s_0+\lambda}-\frac{z}{s_0+z}
=\frac{s_0(\lambda-z)}{(s_0+\lambda)(s_0+z)} ,
\]
hence
\[
Q_{s_0}(\zeta)
=\int_0^\infty\frac{(s_0+\lambda)(s_0+z)}{s_0(\lambda-z)}
\cdot\frac{\dd\sigma(\lambda)}{s_0+\lambda}
=\frac{s_0+z}{s_0}\,q(z),
\]
and $1-\zeta=s_0/(s_0+z)$ gives the first form.

(3) For $z$ near $-s_0$,
\[
\frac1{\lambda-z}=\frac1{(\lambda+s_0)-(z+s_0)}
=\sum_{j\ge0}\frac{(z+s_0)^j}{(\lambda+s_0)^{j+1}},
\]
with uniform convergence on compact subsets of $|z+s_0|<s_0\le
\operatorname{dist}(-s_0,[0,\infty))$; integrating against $\sigma$
identifies $c_j$. Next,
\[
x^k=\Big(1-\frac{s_0}{s_0+\lambda}\Big)^{k}
=\sum_{j=0}^k\binom kj(-s_0)^j\frac1{(s_0+\lambda)^j},
\]
so
\[
\muh_k=\int_0^\infty x^k\,\frac{\dd\sigma(\lambda)}{s_0+\lambda}
=\sum_{j=0}^k\binom kj(-s_0)^j
\int_0^\infty\frac{\dd\sigma(\lambda)}{(s_0+\lambda)^{j+1}}
=\sum_{j=0}^k\binom kj(-s_0)^jc_j .
\]
The matrix of the transform is lower triangular with diagonal entries
$(-s_0)^k\ne0$, hence invertible.
\end{proof}

\begin{remark}
Equation \eqref{eq:binomial} says that the compactification is nothing other
than \emph{re-expansion of the Weyl function about an interior point of the
resolvent set}: compactified moment order $=$ Taylor-jet order at $-s_0$. This
is what makes the moment coordinate intrinsic for fixed $s_0$ --- it requires
no truncation of the underlying spectral data --- and it is also of practical
value: high-precision compactified moments of a measure with intractable
density are obtained from Cauchy-integral Taylor coefficients of $q$ on a
circle about $-s_0$.
\end{remark}

\begin{corollary}[Equal total string length]\label{cor:length}
Suppose in addition $q(0^-)<\infty$. In the Kre\u{\i}n correspondence between
Stieltjes functions and strings, normalized so that the total string length
equals the value of the Weyl function at $z=0$ \cite{KacKrein}, the strings
associated with $q$ and $Q_{s_0}$ have the same total length:
$Q_{s_0}(0)=q(0)$ by \eqref{eq:Qtransform}.
\end{corollary}

\begin{remark}[Non-locality of the bead map]
The equality of lengths is global: the compactification does \emph{not} act
locally on the mass distribution of the string. For a single bead
$[m,\ell]$ one computes $[\,m+1/\ell,\ \ell\,]$, and for two or more beads
the interior masses and gaps mix non-locally (the last mass always maps to
$m_N+1/\ell_N$). We therefore refer to depth in the compactified string as
\emph{resolvent-compactified depth at scale $s_0$}, a bulk coordinate of a
canonically associated equal-length string, not a re-parametrization of the
original one.
\end{remark}

\begin{proposition}[Bounded-transform interpretation; finite vector measures]
\label{prop:operator}
Let $A\ge0$ be self-adjoint on a Hilbert space $\mathcal H$, let
$\phi\in\mathcal H$, and let $\sigma=\sigma_\phi$ be the (finite) spectral
measure of $A$ in $\phi$. Then $B:=A(s_0+A)^{-1}$ is a positive contraction,
$0\le B<1$, and with $\psi:=(s_0+A)^{-1/2}\phi$,
\[
\langle\psi,B^k\psi\rangle=\muh_k\qquad(k\ge0).
\]
\end{proposition}

\begin{proof}
Functional calculus: $\langle\psi,B^k\psi\rangle
=\int\big(\tfrac{\lambda}{s_0+\lambda}\big)^k\tfrac{\dd\sigma_\phi(\lambda)}
{s_0+\lambda}=\muh_k$.
\end{proof}

\begin{remark}[Infinite total mass]\label{rem:infmass}
The measures arising in Section~\ref{sec:zeta} have infinite total mass, so
they are not spectral measures of ordinary Hilbert-space vectors, and
Proposition~\ref{prop:operator} is stated for finite vector measures only.
The main development is measure-theoretic and does not use
Proposition~\ref{prop:operator}. (A rigged-scale formulation is possible,
since \eqref{eq:stieltjescond} says exactly that $(s_0+A)^{-1/2}\phi$ makes
sense for $\phi$ in the natural negative space; we do not pursue this here.)
\end{remark}

% ============================================================================
\section{Finite-rank spectral defects and the first continued-fraction failure}
\label{sec:defects}
% ============================================================================

\subsection{Hankel determinants and the Stieltjes continued fraction}

For a real sequence $(\mu_j)_{j\ge0}$ put
\[
\Delta_n=\det(\mu_{i+j})_{i,j=0}^{n-1},
\qquad
\Delta_n'=\det(\mu_{i+j+1})_{i,j=0}^{n-1},
\qquad
\Delta_0=\Delta_0'=1 .
\]
We use throughout the string form of the Stieltjes continued fraction,
\begin{equation}\label{eq:cf}
q(z)=
\cfrac{1}{-zm_1+\cfrac{1}{\ell_1+\cfrac{1}{-zm_2+\cfrac{1}{\ell_2+\cdots}}}} ,
\end{equation}
associated with $q(z)=\int(\,t-z)^{-1}\dd\mu(t)=-\sum_{k\ge0}\mu_kz^{-k-1}$
as a formal expansion at $z=\infty$; in the Kre\u{\i}n string picture the
$m_k$ are point masses and the $\ell_k$ the gaps between them, with the first
mass at the measuring end \cite{KacKrein,GK}. The classical determinant
formulas, in this convention, are
\begin{equation}\label{eq:detformulas}
m_k=\frac{(\Delta'_{k-1})^2}{\Delta_{k-1}\Delta_k},
\qquad
\ell_k=\frac{\Delta_k^2}{\Delta'_{k-1}\Delta'_k}
\qquad(k\ge1),
\end{equation}
see \cite{Stieltjes,Wall,GK}; we have re-verified
\eqref{eq:detformulas} in this exact indexing symbolically for one and two
atoms and numerically to $10^{-112}$ for generic six-atom measures.
Positivity of all $m_k,\ell_k$ up to level $N$ is equivalent to
$\Delta_n>0$ and $\Delta_n'>0$ up to the corresponding orders --- the
classical solvability criterion for the Stieltjes moment problem
\cite{Stieltjes,Schmudgen}.

\begin{theorem}[Exact index of the first continued-fraction failure]
\label{thm:cfindex}
Let $(\mu_j)$ be a real sequence whose Hankel determinants $\Delta_n$,
$\Delta_n'$ are nonzero up to the orders below. Define
\[
n_0=\min\{n\ge1:\ \Delta_{n-1}\Delta_n<0\},
\qquad
n_1=\min\{n\ge1:\ \Delta'_{n-1}\Delta'_n<0\}
\]
(with $\min\emptyset=\infty$). If $\min(n_0,n_1)<\infty$, then all
coefficients of \eqref{eq:cf} preceding the following one are positive, and
the first nonpositive coefficient is
\[
\begin{cases}
m_{n_0}, & \text{raw index } 2n_0-2,\quad\text{if } n_0\le n_1,\\[2pt]
\ell_{n_1}, & \text{raw index } 2n_1-1,\quad\text{if } n_1<n_0 .
\end{cases}
\]
(Coefficients are indexed $m_1,\ell_1,m_2,\ell_2,\dots$ with raw indices
$0,1,2,3,\dots$; a tie $n_0=n_1$ appears first as a mass since $2n-2<2n-1$.)
\end{theorem}

\begin{proof}
By \eqref{eq:detformulas} the numerators are squares, so
$\sgn m_k=\sgn(\Delta_{k-1}\Delta_k)$ and
$\sgn\ell_k=\sgn(\Delta'_{k-1}\Delta'_k)$. All factors before the first sign
change in each family are positive; the first negative product in either
family occurs exactly at $n_0$, resp.\ $n_1$, and the interleaving
$m_1,\ell_1,m_2,\ell_2,\dots$ places $m_{n_0}$ at raw index $2n_0-2$ and
$\ell_{n_1}$ at $2n_1-1$.
\end{proof}

\begin{remark}
Theorem~\ref{thm:cfindex} is the bridge between moment-space detection and
inverse-string depth: \emph{the first anomalous continued-fraction
coefficient is exactly the first order at which the truncated Stieltjes
moment problem detects the defect}, and the parity of the anomaly (mass
versus length) records which Hankel family failed first. Both statements are
convention-checked facts about \eqref{eq:detformulas}; the content of the
following two theorems is \emph{where} that failure occurs.
\end{remark}

\subsection{Rank-one defects: a forbidden real atom}

\begin{theorem}[Exact rank-one Christoffel threshold]\label{thm:rankone}
Let $\mu_0$ be a finite positive measure with compact support
$K\subset[0,\infty)$, with Hankel matrices $H_n=(\mu_{i+j})_{i,j=0}^{n-1}$
and $H_n'=(\mu_{i+j+1})_{i,j=0}^{n-1}$ nonsingular, and let
\[
\mu_\eps=\mu_0+\eps\,\delta_a,\qquad a<0,\ \eps>0 .
\]
Then all ordinary Hankel determinants of $\mu_\eps$ remain positive, while
\begin{equation}\label{eq:rankone}
\frac{\det H_n'(\eps)}{\det H_n'(0)}
=1+\eps\,a\,K'_{n-1}(a,a),
\end{equation}
where $K'_{n-1}(a,a)=v(a)^{T}(H_n')^{-1}v(a)$,
$v(a)=(1,a,\dots,a^{n-1})^{T}$, is the (diagonal) Christoffel--Darboux kernel
of the positive measure $t\dd\mu_0(t)$. Consequently the first Stieltjes
failure is a shifted-Hankel failure, at the exact threshold
\begin{equation}\label{eq:threshold}
n_{\det}=\min\{\,n:\ \eps\,|a|\,K'_{n-1}(a,a)>1\,\},
\end{equation}
and by Theorem~\ref{thm:cfindex} it appears as the length coefficient
$\ell_{n_{\det}}$, at raw continued-fraction index $2n_{\det}-1$.
\end{theorem}

\begin{proof}
$\mu_\eps$ is a positive measure on $\R$, so its ordinary Hankel matrices are
positive semidefinite Gram matrices; with $H_n(0)>0$ and $\eps>0$ they are
positive definite. For the shifted family, the moment increments are
$\delta\mu_{k}=\eps a^{k}$, so
$H_n'(\eps)=H_n'(0)+\eps a\,v(a)v(a)^{T}$ is a rank-one update, and the
matrix determinant lemma gives
$\det H_n'(\eps)=\det H_n'(0)\,\big(1+\eps a\,v(a)^{T}(H_n')^{-1}v(a)\big)$.
The identification of $v^{T}(H')^{-1}v$ with the CD kernel of $t\dd\mu_0$ is
the classical Christoffel-function identity: if $(p_j')$ are the orthonormal
polynomials of $t\dd\mu_0$, then $(H_n')^{-1}=P^{T}P$ where $P$ maps
coefficient vectors to $(p'_j)$-coordinates, so
$v^{T}(H_n')^{-1}v=\sum_{j\le n-1}p_j'(a)^2=K'_{n-1}(a,a)$
\cite{Nevai,Totik}. Since $a<0$ and $K'_{n-1}(a,a)>0$ is nondecreasing in
$n$, the ratio \eqref{eq:rankone} decreases through $0$ exactly at
\eqref{eq:threshold}; the ordinary family never fails, so
Theorem~\ref{thm:cfindex} applies with $n_1=n_{\det}<n_0=\infty$.
\end{proof}

\subsection{Rank-two defects: a conjugate non-real pair}

\begin{theorem}[Exact rank-two determinant identity]\label{thm:ranktwo}
Let $\mu_0$ be as above, $c\in\C\setminus\R$, and consider the real moment
perturbation
\[
\mu_k(\eps)=\mu_k+\frac\eps2\,c^k+\frac\eps2\,\bar c^{\,k}\qquad(k\ge0),
\]
i.e.\ formal conjugate atoms $\tfrac\eps2\delta_c+\tfrac\eps2\delta_{\bar c}$.
Define the bilinear kernel $K^{\mathrm{bil}}_{n-1}(z,w)=v(z)^{T}H_n^{-1}v(w)
=\sum_{j\le n-1}p_j(z)p_j(w)$ and the Hermitian value
$K_{n-1}(c,\bar c)=\sum_{j\le n-1}|p_j(c)|^2$, where $(p_j)$ are the
orthonormal polynomials of $\mu_0$. Then
\begin{equation}\label{eq:ranktwo}
\frac{\det H_n(\eps)}{\det H_n(0)}
=1+\eps\,\Rea K^{\mathrm{bil}}_{n-1}(c,c)
+\frac{\eps^2}4\Big(\big|K^{\mathrm{bil}}_{n-1}(c,c)\big|^2
-K_{n-1}(c,\bar c)^2\Big),
\end{equation}
and the coefficient of $\eps^2/4$ is nonpositive, by Cauchy--Schwarz
$\big|\sum p_j(c)^2\big|\le\sum|p_j(c)|^2$. The analogous formula holds for
the shifted family $H_n'$, with kernels of $t\dd\mu_0$ and defect columns
weighted by $c$, $\bar c$.
\end{theorem}

\begin{proof}
$H_n(\eps)=H_n+U C U^{*T}$ in the bilinear sense, with
$U=[v(c)\ v(\bar c)]$ and $C=\operatorname{diag}(\eps/2,\eps/2)$; since the
perturbation is symmetric (not Hermitian) in the Hankel structure, apply the
two-column matrix determinant lemma
$\det(H+U C V^{T})=\det H\cdot\det(I_2+C V^{T}H^{-1}U)$ with
$V=[v(\bar c)\ v(c)]$, so that the moment increments
$\tfrac\eps2(c^{i+j}+\bar c^{\,i+j})$ are reproduced. Writing
$V^{T}H^{-1}U=\begin{pmatrix}K^{\mathrm{bil}}(\bar c,c)&
K^{\mathrm{bil}}(\bar c,\bar c)\\K^{\mathrm{bil}}(c,c)&
K^{\mathrm{bil}}(c,\bar c)\end{pmatrix}$ and using that the background
moments are real, so that
$K^{\mathrm{bil}}(\bar c,\bar c)=\overline{K^{\mathrm{bil}}(c,c)}$ and
$K^{\mathrm{bil}}(c,\bar c)=K^{\mathrm{bil}}(\bar c,c)
=K_{n-1}(c,\bar c)\in(0,\infty)$, the $2\times2$ determinant expands to
\eqref{eq:ranktwo}. Cauchy--Schwarz in $\ell^2$ gives the sign of the
quadratic term.
\end{proof}

\begin{remark}[No factor-of-two shortcut]\label{rem:nofactor2}
The negative channel in \eqref{eq:ranktwo} grows on the exponential scale
like $\eps^2e^{4ng(c)}$, but the defect weight is squared: solving either
threshold gives the same leading order $n\sim\log(1/\eps)/(2g(c))$ at fixed
defect location. Observed depth differences between defect types in our
laboratory experiments came from different Green distances of the defect
locations, not from the exponent $4g$.
\end{remark}

\begin{remark}[Finitely many defects]\label{rem:multidefect}
For $r$ defects at $c_1,\dots,c_r$ with weight matrix $C$, the same lemma
gives $\det H_n(\eps)/\det H_n=\det(I_r+C\,\mathcal K_n)$ with
$\mathcal K_n$ the $r\times r$ matrix of CD-kernel evaluations at the defect
points: an exact finite-dimensional interaction law. On the exponential
scale the first failure is governed by the defect(s) of smallest Green
distance; we leave the multi-defect asymptotics to future work.
\end{remark}

% ============================================================================
\section{Christoffel detection depth}\label{sec:depth}
% ============================================================================

Let $K\subset\R$ be compact, non-polar, with Green function $g_K$ of
$\overline\C\setminus K$ with pole at infinity, and let $\mu_0$ be a positive
measure with $\supp\mu_0=K$ that is \emph{regular} in the Stahl--Totik sense
\cite{StahlTotik}, i.e.\ its orthonormal polynomials obey the $n$-th root
asymptotics $|p_n(z)|^{1/n}\to e^{g_K(z)}$ locally uniformly on
$\C\setminus\operatorname{Co}(K)$, and hence
\begin{equation}\label{eq:rootasym}
\frac1{2n}\log K_n(c,\bar c)\longrightarrow g_K(c)
\qquad(n\to\infty)
\end{equation}
for fixed $c\notin K$, where $K_n(c,\bar c)=\sum_{j\le n}|p_j(c)|^2$.
(For real $a\notin\operatorname{Co}(K)$ the same holds for $K_n(a,a)$; we
also apply \eqref{eq:rootasym} to the measure $t\dd\mu_0$, which is regular
if $\mu_0$ is and $0\notin K$.)

\begin{theorem}[Fixed-defect detection depth]\label{thm:depth}
Let $\mu_0$ be regular with compact non-polar support
$K\subset(0,\infty)$.
\begin{enumerate}
\item (Rank one.) For a defect $\eps\delta_a$ with fixed
$a<0$, the first failure order $n_\eps$ of \eqref{eq:threshold} satisfies
\[
\frac{n_\eps}{\log(1/\eps)}\longrightarrow\frac1{2g_K(a)}
\qquad(\eps\downarrow0).
\]
\item (Rank two.) For a conjugate defect pair at fixed
$c\in\C\setminus\R$ as in Theorem~\ref{thm:ranktwo}, assume the
nondegeneracy condition that the negative channel obeys the root asymptotic
$\frac1{4n}\log\big(K_{n-1}(c,\bar c)^2-
|K^{\mathrm{bil}}_{n-1}(c,c)|^2\big)\to g_K(c)$. Then the first failure
order of the relevant Hankel family satisfies
\[
\frac{n_\eps}{\log(1/\eps)}\longrightarrow\frac1{2g_K(c)}
\qquad(\eps\downarrow0).
\]
\end{enumerate}
\end{theorem}

\begin{proof}
(1) Fix $\eps_0>0$. By \eqref{eq:rootasym} applied to $t\dd\mu_0$ (also
regular, with the same exterior Green function), and since the fixed factor
$|a|>0$ is absorbed by any exponent tolerance, there is $n_*$ such that
\[
e^{2n(g-\eps_0)}\ \le\ |a|\,K'_{n-1}(a,a)\ \le\ e^{2n(g+\eps_0)}
\qquad(n\ge n_*).
\]
The threshold \eqref{eq:threshold} is crossed at the first $n$ with
$\eps\,|a|K'_{n-1}(a,a)>1$; since $K'_n$ is nondecreasing in $n$, the
two-sided bound yields, for all small $\eps$,
\[
\frac{\log(1/\eps)}{2(g+\eps_0)}-O(1)
\le n_\eps\le
\frac{\log(1/\eps)}{2(g-\eps_0)}+O(1),
\]
and $\eps_0\downarrow0$ gives the claim.

(2) Write \eqref{eq:ranktwo} as $1+\eps L_n-\tfrac{\eps^2}4 N_n$ with
$L_n=\Rea K^{\mathrm{bil}}_{n-1}(c,c)$, so $|L_n|\le K_{n-1}(c,\bar c)$, and
$N_n=K_{n-1}(c,\bar c)^2-|K^{\mathrm{bil}}_{n-1}(c,c)|^2\ge0$. Write
$L:=\log(1/\eps)$. \emph{No failure below the window:} for
$n\le(1-\kappa)L/(2g)$ and $\eps_0<g\kappa/(1-\kappa)$, both
$\eps|L_n|\le\eps\,e^{2n(g+\eps_0)}\to0$ and
$\eps^2N_n\le\eps^2e^{4n(g+\eps_0)}\to0$, so the ratio stays positive.
\emph{Failure above the window:} for $n\ge(1+\kappa)L/(2g)$ the
nondegeneracy hypothesis gives, for $\eps_0$ small,
\[
\tfrac{\eps^2}4N_n\ \ge\ \tfrac14\,e^{-2L+4n(g-\eps_0)}
\ \ge\ \tfrac14\,e^{2\kappa'L},
\qquad
\eps|L_n|\ \le\ e^{-L+2n(g+\eps_0)}\ \le\ e^{(\kappa'+o(1))L}
\]
with $\kappa'>0$; hence
$\eps^2N_n/(\eps|L_n|)\ \ge\ \tfrac14 e^{(\kappa'-o(1))L}\to\infty$, and the
determinant ratio $1+\eps L_n-\tfrac{\eps^2}4N_n$ is negative for all small
$\eps$, whatever the sign of $L_n$. Since a monotone-in-$n$ failure witness
is not needed --- we only require that no $n$ below the window fails and
some $n$ above it fails --- the first failure order lies in
$[(1-\kappa),(1+\kappa)]\cdot L/(2g)$, and $\kappa\downarrow0$ completes the
proof.
\end{proof}

\begin{remark}[What regularity does and does not give]\label{rem:noO1}
Stahl--Totik regularity supplies $n$-th root asymptotics and hence the ratio
limit of Theorem~\ref{thm:depth}. It does \emph{not} supply the
multiplicative constants that would be needed for a statement of the form
$n_\eps=\log(C/\eps)/(2g)+O(1)$; such refinements require strong (Szeg\H o
class) asymptotics and additional hypotheses. Our numerical experiments
confirm this caution: the effective constant varies substantially with the
background and with the compactification scale (Section~\ref{sec:numerics}).
\end{remark}

\begin{corollary}[Continued-fraction defect depth]\label{cor:cfdepth}
In the setting of Theorem~\ref{thm:depth}, the first anomalous
continued-fraction coefficient of the defected system occurs at raw index
$2n_\eps+O(1)$ with the mass/length parity of Theorem~\ref{thm:cfindex}:
Green distance from the admissible support is, asymptotically, the
inverse-spectral resolution depth.
\end{corollary}

% ============================================================================
\section{The shifted zeta-string reference family}\label{sec:zeta}
% ============================================================================

\subsection{The positive reference system}
Let $\xi(s)=\tfrac12 s(s-1)\pi^{-s/2}\Gamma(s/2)\zeta(s)$ and, for
$\omega>0$, define $q_\omega$ by \eqref{eq:qomega}. The following facts are
established in \cite{Suzuki2023,MarginPaper,StringPaper}: if the half-plane
$\Rea s>\tfrac12+\omega$ is free of zeros of $\xi$ --- in particular, for
every $\omega>0$, under the Riemann hypothesis --- then $q_\omega$ is a
Stieltjes function; its spectral measure $\sigma_\omega$ is absolutely
continuous on $(0,\infty)$ with density
\begin{equation}\label{eq:density}
\rho_\omega(\lambda)
=\frac{1}{\pi\sqrt{\lambda}}\,
\Rea\frac{\xi'}{\xi}\Big(\tfrac12+\omega-i\sqrt\lambda\Big)
=\frac{1}{\pi\sqrt\lambda}\sum_{\gamma}
\frac{\omega}{\omega^2+(\sqrt\lambda-\gamma)^2},
\end{equation}
the sum running over the ordinates of the (paired) zeros: the
Poisson-broadened zero measure. The density is continuous and strictly
positive on $(0,\infty)$, with $\rho_\omega(\lambda)\sim
S_\omega/(\pi\sqrt\lambda)$ as $\lambda\downarrow0$, where
$S_\omega=(\xi'/\xi)(\tfrac12+\omega)>0$, and
$\rho_\omega(\lambda)\asymp\log\lambda/\sqrt\lambda$ as
$\lambda\to\infty$; in particular
$\int(1+\lambda)^{-1}\dd\sigma_\omega<\infty$, so
Theorem~\ref{thm:compactification} applies. We call
$(\sigma_\omega,q_\omega)$ the \emph{positive shifted-zeta reference
system}. After compactification at scale $s_0$, the reference measure
$\mu_{\omega,s_0}$ is a finite measure on $[0,1]$ whose density is
continuous and positive on $(0,1)$; in particular $\mu_{\omega,s_0}$ is
regular in the Stahl--Totik sense \cite{StahlTotik}.

\subsection{One hypothetical quartet}

\begin{proposition}[Location and weight of the defect]\label{prop:quartet}
Insert into \eqref{eq:qomega} a hypothetical simple zero
$\rho=\tfrac12+\delta+i\tau$ with $\delta>\omega>0$, $\tau>0$, together with
the conjugate and functional-equation partners
$\bar\rho,\,1-\rho,\,1-\bar\rho$. Set $\eta=\delta-\omega>0$. On the
principal branch $\Rea w>0$, the added quartet contributes to $q_\omega$
exactly two poles, located at
\[
z_\pm=-(\eta\pm i\tau)^2=\tau^2-\eta^2\mp 2i\eta\tau ,
\]
each with Stieltjes-form weight $2$:
$q_\omega(z)\sim 2/(z_\pm-z)$ as $z\to z_\pm$. Under resolvent
compactification at scale $s_0$ the defect points and weights become
\[
c_\pm=\frac{z_\pm}{s_0+z_\pm},
\qquad
\eps_{\mathrm{eff},\pm}=\frac{2}{s_0+z_\pm},
\qquad
|\eps_{\mathrm{eff}}|\sim\frac{2}{\tau^2}\ \ (\tau\to\infty).
\]
The partners $1-\rho,1-\bar\rho$ produce no principal-sheet poles; together
with the cut contribution of the quartet they produce a signed absolutely
continuous correction treated in Lemma~\ref{lem:cut}.
\end{proposition}

\begin{proof}
$\xi'/\xi(\tfrac12+\omega+w)$ has a simple pole of residue $1$ at each zero;
the zero $\rho$ gives $w_0=\rho-\tfrac12-\omega=\eta+i\tau$ with
$\Rea w_0>0$ (principal sheet), and $\bar\rho$ gives $\bar w_0$; the
partners give $\Rea w<0$. Near $w_0$, with $z=-w^2$,
$z-z_+=-2w_0(w-w_0)+O((w-w_0)^2)$, so
\[
q_\omega(z)=\frac{1}{w}\cdot\frac{1}{w-w_0}(1+o(1))
=\frac{1}{w_0}\cdot\frac{-2w_0}{z-z_+}(1+o(1))
=\frac{2}{z_+-z}(1+o(1)).
\]
The compactified data follow from Theorem~\ref{thm:compactification}: a pole
$w/(\lambda_0-z)$ transforms to a pole at $x_0=\lambda_0/(s_0+\lambda_0)$
with weight $w/(s_0+\lambda_0)$, applied at $\lambda_0=z_\pm$.
\end{proof}

\subsection{The signed cut correction}

Beyond its two principal poles, the quartet alters the boundary values of
$q_\omega$ on the cut $[0,\infty)$: writing the quartet's exact contribution
to $(\xi'/\xi)(\tfrac12+\omega+w)$ as
$d(w)=\sum_{\gamma_v\in\{\tau-i\delta,\ \tau+i\delta\}}
\big(\tfrac1{\,\omega+w-i\gamma_v}+\tfrac1{\,\omega+w+i\gamma_v}\big)$,
the function $q_{\mathrm{def}}(z)=d(\sqrt{-z})/\sqrt{-z}$ decomposes as the
two poles of Proposition~\ref{prop:quartet} plus an absolutely continuous
part on $[0,\infty)$ with signed density
\begin{equation}\label{eq:cutdensity}
\rho_{\mathrm{def}}(\lambda)
=\frac{1}{\pi\sqrt{\lambda}}\,\Rea d\big(-i\sqrt\lambda\,\big) .
\end{equation}
When the defect is scaled by a small parameter $\eps$ (the sensitivity
model), the correction is $\eps\rho_{\mathrm{def}}$.

\begin{lemma}[The cut correction cannot affect the leading depth]
\label{lem:cut}
Fix $\omega>0$, $s_0>0$ and the quartet parameters $(\tau,\delta)$,
$\delta>\omega$. Then
\[
M_*:=\sup_{\lambda>0}
\frac{|\rho_{\mathrm{def}}(\lambda)|}{\rho_\omega(\lambda)}<\infty .
\]
Consequently, for $0<\eps<1/(2M_*)$, the perturbed background
$\dd\sigma_\omega+\eps\rho_{\mathrm{def}}\dd\lambda$ satisfies, in quadratic
form on every polynomial subspace and for every $n$,
\[
(1-\eps M_*)\,H_n^{(0)}\ \le\ H_n^{(0)}+\eps R_n\ \le\ (1+\eps M_*)\,H_n^{(0)},
\]
where $H_n^{(0)}$ and $R_n$ are the compactified Hankel matrices of
$\sigma_\omega$ and of the cut correction. Hence the measure
``background $+$ cut correction'' remains a positive measure with the same
support, regularity class, and Green function, its Christoffel kernels are
distorted by factors in $[(1+\eps M_*)^{-1},(1-\eps M_*)^{-1}]$, and the
first failure order of the full defected system equals that of the
principal-pole (rank-two) model up to an $O(1)$ shift as $\eps\downarrow0$:
\[
n^{\mathrm{full}}_\eps=n^{\mathrm{pole}}_\eps+O(1)
\qquad(\eps\downarrow0).
\]
In particular the ratio limit of Theorem~\ref{thm:depth} is unchanged.
\end{lemma}

\begin{proof}
\emph{Boundedness of the ratio.} Both densities are continuous and positive
(resp.\ continuous) on $(0,\infty)$, so it suffices to compare the ends. As
$\lambda\downarrow0$: $\Rea d(-i\sqrt\lambda)$ tends to the finite constant
$d(0)$, while $\Rea(\xi'/\xi)(\tfrac12+\omega-i\sqrt\lambda)\to S_\omega>0$;
both densities behave like $\text{const}/\sqrt\lambda$ and the ratio tends to
$|d(0)|/S_\omega$. As $\lambda\to\infty$: with $u=\sqrt\lambda$,
$\Rea d(-iu)=O(u^{-2})$ (two Lorentzian pairs centred at $\pm\tau$), while
$\Rea(\xi'/\xi)(\tfrac12+\omega-iu)\asymp\log u\to\infty$ by
\eqref{eq:density} and the Riemann--von Mangoldt density; the ratio tends to
$0$. Hence $M_*<\infty$.

\emph{Quadratic-form domination.} For any polynomial $P$ (in the
compactified variable; the compactification map is a fixed smooth bijection,
so pointwise domination of densities is preserved),
\[
\Big|\int P^2\,\eps\rho_{\mathrm{def}}\Big|
\le\eps M_*\int P^2\,\rho_\omega
=\eps M_*\,\langle P,H^{(0)}P\rangle .
\]
\emph{Threshold stability.} Let $\widetilde H_n=H^{(0)}_n+\eps R_n$; by the
two-sided bound, $\widetilde H_n$ is positive definite for
$\eps<1/M_*$ and its inverse satisfies
$(1+\eps M_*)^{-1}(H^{(0)}_n)^{-1}\le\widetilde H_n^{-1}\le
(1-\eps M_*)^{-1}(H^{(0)}_n)^{-1}$, hence the CD kernels obey
$\widetilde K_n(c,\bar c)\in
[(1+\eps M_*)^{-1},(1-\eps M_*)^{-1}]\cdot K_n(c,\bar c)$,
and the bilinear values are controlled by the same two-sided operator bound
via polarization, uniformly in $n$, once $\eps\le 1/(2M_*)$. The exact
rank-two threshold of Theorem~\ref{thm:ranktwo}, computed against
$\widetilde H_n$ instead of $H^{(0)}_n$, is therefore displaced by at most a
fixed multiplicative factor $1+O(\eps)$ inside the logarithm, i.e.\ the first
failure order shifts by $O(\eps)/(2g)+O(1)=O(1)$.
\end{proof}

\begin{remark}
Lemma~\ref{lem:cut} discharges what \cite{ROUND6} flagged as the one
remaining local proof obligation: the finite-rank model is not an
approximation but the exact leading mechanism, the a.c.\ correction being
dominated by the reference density. The same argument applies verbatim to
any finite collection of quartets.
\end{remark}

\subsection{The sensitivity theorem}

We emphasize the logical status of the following statement: it concerns the
positive reference system \eqref{eq:qomega}--\eqref{eq:density} perturbed by
\emph{one hypothetical quartet scaled by $\eps$}. It is a perturbative
sensitivity theorem about an exactly defined model. It is \emph{not} a
statement about the actual zeta function in a world where RH fails: there,
nothing would confine the violation to a single quartet, and we make no such
claim. (A conditional variant --- assuming all remaining zeros admissible,
the actual response decomposes as reference plus the quartet's exact
contribution, to which the same analysis applies by
Lemma~\ref{lem:cut} --- holds with identical proof.)

\begin{theorem}[Spectral-defect sensitivity of the shifted-zeta reference
system]\label{thm:zetadepth}
Fix $\omega>0$, $s_0>0$, and quartet parameters $\tau>0$,
$\eta=\delta-\omega>0$. Let $\mu_{\omega,s_0}$ be the compactified reference
measure, $g_{\omega,s_0}$ the Green function of the complement of its
support, and $c_\pm$ the compactified defect points of
Proposition~\ref{prop:quartet}. Scale the quartet by $\eps\downarrow0$ and
let $n_\eps$ be the first compactified Hankel order at which the defected
system loses Stieltjes positivity. Then, under the nondegeneracy condition
of Theorem~\ref{thm:depth}(2) at $c_+$,
\[
\frac{n_\eps}{\log(1/\eps)}
\longrightarrow
\frac1{2\,g_{\omega,s_0}(c_+)},
\]
and the first anomalous continued-fraction coefficient of the compactified
string occurs at the index given by Theorem~\ref{thm:cfindex}.
\end{theorem}

\begin{proof}
The compactified reference measure is finite with support
$\overline{(0,1)}=[0,1]$ and positive continuous density on $(0,1)$, hence
regular \cite{StahlTotik}; Proposition~\ref{prop:quartet} realizes the
scaled quartet as the conjugate rank-two defect
$\tfrac{\eps'}2\delta_{c_+}+\tfrac{\bar\eps'}2\delta_{c_-}$ with
$\eps'=2\eps/(s_0+z_+)$, plus the cut correction; Lemma~\ref{lem:cut}
reduces to the pure finite-rank model; Theorems~\ref{thm:ranktwo}
and~\ref{thm:depth}(2) give the ratio limit; Theorem~\ref{thm:cfindex}
locates the continued-fraction anomaly.
\end{proof}

\subsection{High height: a moving-point asymptotic}\label{sec:highheight}

For $\tau\to\infty$ at fixed $\eta,s_0,\omega$, the defect point approaches
the compactified endpoint:
\[
1-\Rea c_+\sim\frac{s_0}{\tau^2},
\qquad
|\Ima c_+|\sim\frac{2s_0\eta}{\tau^3},
\]
and the equilibrium density of the limiting interval near its upper edge
gives the explicit Green asymptotic
\begin{equation}\label{eq:gmoving}
g(c_+)\sim\frac{2\eta\sqrt{s_0}}{\tau^2}
\qquad\Big(\text{more precisely }
2\eta\sqrt{s_0+\lambda_{\min}}/\tau^2\text{ for a support hull
}[\lambda_{\min},\infty)\Big),
\end{equation}
verified against the exact interval Green function to $0.4\%$ at
$\tau=160$ \cite{LabRecord}. Combining \eqref{eq:gmoving} with
$|\eps_{\mathrm{eff}}|\sim2/\tau^2$ \emph{formally} in the fixed-point
threshold yields the high-height prediction
\begin{equation}\label{eq:highheight}
n_{\mathrm{comp}}
\;\asymp\;
\frac{\tau^2\log\tau}{2\eta\sqrt{s_0}} .
\end{equation}
We state \eqref{eq:highheight} as a \emph{moving-point asymptotic supported
numerically}: both the evaluation point and the effective weight vary with
$\tau$, and fixed-point Stahl--Totik regularity does not supply the required
uniformity. Establishing uniform Christoffel asymptotics in the joint regime
$n\asymp\tau^2\log\tau/(\eta\sqrt{s_0})$ --- by strong asymptotics for the
compactified $\omega$-density, endpoint asymptotics plus analytic
continuation, Riemann--Hilbert methods, or extremal polynomial bounds --- is
an open problem we consider tractable but do not resolve here. The numerical
section documents both the regime where \eqref{eq:highheight}-type
predictions hold (broadening or displacement exceeding the local zero
spacing) and the quasi-discrete regime of genuine low-height zeta data,
where the interval law overestimates the observed depth.

\subsection{Dependence on the scale $s_0$}\label{sec:s0}

The resolvent compactification is a one-parameter family, and the detection
order genuinely depends on the scale: in controlled experiments with fixed
spectral data and defect, $n_{s_0}$ is non-monotone in $s_0$ (e.g.\
$22,14,12,17,28$ at $s_0=1,4,25,100,400$ for a fixed moderate-height defect
\cite{LabRecord}), and no power-law renormalization removes the dependence
at moderate scales. At high height with $\tau^2\gg s_0$, \eqref{eq:gmoving}
gives the leading covariance $n_{s_0}\propto1/\sqrt{s_0+\lambda_{\min}}$.
We therefore use throughout the terminology \emph{resolvent-compactified
detection depth at scale $s_0$}: it is cutoff-independent for fixed $s_0$
(Theorem~\ref{thm:compactification} makes it a Taylor-jet order), but it is
not scale-free. The observed minimum of $n_{s_0}$ near $s_0\sim\tau^2$
suggests a detection-optimal scale $s_0^*(\tau)$; we leave this as a
future-work observation.

% ============================================================================
\section{Time-domain detection}\label{sec:time}
% ============================================================================

We now compare with the shifted screw-function criterion of
\cite{Suzuki2023,MarginPaper}. For fixed $\omega>0$ the positive reference
background has the linear margin
\[
\Psi^{\mathrm{bg}}_\omega(t)=S_\omega\,t+O(1),
\qquad
S_\omega=\frac{\xi'}{\xi}\Big(\tfrac12+\omega\Big)>0 ,
\]
and one violating quartet contributes the growing oscillation
\begin{equation}\label{eq:quartetgrow}
\Psi^{\mathrm{grow}}_{\omega,\rho}(t)
=-\frac{2e^{\eta t}}{\tau^2+\eta^2}\cos(\tau t+\alpha)+O(t),
\end{equation}
with amplitude coefficient $2/(\tau^2+\eta^2)$; this normalization has been
audited symbolically and numerically against the exact spectral series (the
two growing branches are $e^{(\eta\pm i\tau)t}/(\eta\pm i\tau)^2$;
agreement of \eqref{eq:quartetgrow} with the exact quartet term to relative
$10^{-8}$ by $t=150$) \cite{LabRecord,MarginPaper}.

\begin{theorem}[Envelope detection time]\label{thm:envelope}
Define $t_{\mathrm{env}}$ by
$\dfrac{2e^{\eta t_{\mathrm{env}}}}{\tau^2+\eta^2}=S_\omega\,
t_{\mathrm{env}}$. Then, for fixed $\eta,\omega>0$, as $\tau\to\infty$,
\[
\eta\,t_{\mathrm{env}}
=2\log\tau+\log\log\tau+\log\frac{S_\omega}{\eta}+o(1).
\]
\end{theorem}

\begin{proof}
Set $x=\eta t_{\mathrm{env}}$; the defining equation reads
$2e^x/(\tau^2+\eta^2)=(S_\omega/\eta)\,x$, i.e.
\[
x=\log(\tau^2+\eta^2)+\log x+\log\frac{S_\omega}{2\eta}.
\]
Since $\log(\tau^2+\eta^2)=2\log\tau+o(1)$ and a first iteration gives
$x=2\log\tau+O(\log\log\tau)$, we get
$\log x=\log\log\tau+\log2+o(1)$, and substituting back,
$x=2\log\tau+\log\log\tau+\log2+\log\frac{S_\omega}{2\eta}+o(1)
=2\log\tau+\log\log\tau+\log\frac{S_\omega}{\eta}+o(1)$.
\end{proof}

\begin{theorem}[First-negative-time bracket]\label{thm:bracket}
Let $\Psi(t)=S_\omega t+\Psi_{\omega,\rho}(t)$ be the one-quartet
perturbative model with the exact quartet term, and
$t_{\mathrm{first}}=\inf\{t>0:\Psi(t)<0\}$. Assume the transversality
condition that at $t_{\mathrm{env}}$ the envelope crosses the (exact,
non-oscillatory) margin at the rate $\eta+o(1)$, and that the
non-oscillatory part varies by $o(1)$ relative to its value on intervals of
length $O(1/\tau)$. Then, for $\tau$ large,
\[
0\ \le\ t_{\mathrm{first}}-t_{\mathrm{env}}\ \le\ \frac{2\pi}\tau
+o(1/\tau),
\qquad\text{hence}\qquad
\eta\,t_{\mathrm{first}}
=2\log\tau+\log\log\tau+\log\frac{S_\omega}{\eta}+o(1).
\]
\end{theorem}

\begin{proof}
Before $t_{\mathrm{env}}$ the envelope is below the margin, so
$\Psi>0$: $t_{\mathrm{first}}\ge t_{\mathrm{env}}$. Within
$(t_{\mathrm{env}},t_{\mathrm{env}}+2\pi/\tau]$ the cosine attains a trough
$\cos=-1$; over that interval the envelope changes by the factor
$e^{2\pi\eta/\tau}=1+O(1/\tau)$ and the non-oscillatory part by $o(1)$
relative terms, so at the trough
$\Psi\le S_\omega t\,(1+o(1))-\text{envelope}(t)\,(1-O(1/\tau))<0$ by
transversality. Numerically the bracket is comfortable: gaps
$0.115,0.027,0.011,0.0035$ against bounds $2\pi/\tau=
0.314,0.079,0.020,0.0049$ at $\tau=20,80,320,1280$ \cite{LabRecord}.
\end{proof}

\begin{corollary}[The two coordinates compared]\label{cor:compare}
For the same hypothetical quartet at height $\tau$ with horizontal excess
$\eta=\delta-\omega$: the time-domain criterion detects it at
\[
t_{\mathrm{first}}=\frac{2\log\tau+\log\log\tau+O(1)}{\eta},
\]
by Theorems \ref{thm:envelope}--\ref{thm:bracket}, while the
resolvent-compactified inverse-spectral coordinate at scale $s_0$ requires
depth $n_\eps\sim\log(1/\eps)/(2g_{\omega,s_0}(c_+))$ for the scaled defect
(Theorem~\ref{thm:zetadepth}), with the high-height moving-point behavior
\eqref{eq:highheight} predicting $n\asymp\tau^2\log\tau/(2\eta\sqrt{s_0})$.
The time-domain coordinate resolves a high defect on a logarithmic-in-height
scale; the compactified moment description is predicted, and numerically
observed in the accessible regimes, to require polynomially-in-height larger
resolution. This is a statement about the conditioning of two exact
representations of the same object, in the stated coordinates and asymptotic
regimes.
\end{corollary}

% ============================================================================
\section{Numerical illustrations}\label{sec:numerics}
% ============================================================================

All computations below are reproducible from the laboratory record
\cite{LabRecord}; multi-precision arithmetic (60--450 digits), exact
integer certification where applicable, and control rows (undefected
backgrounds remaining Stieltjes-clean) accompany every experiment. Numerics
here illustrate regimes and constants; no proof depends on them.

\subsection{Exactness certificates}
The binomial transform \eqref{eq:binomial} was verified on rational toy
measures to $8\cdot10^{-131}$; the rank-one identity \eqref{eq:rankone} to
$10^{-285}$; the rank-two identity \eqref{eq:ranktwo} to $10^{-273}$; the
determinant formulas \eqref{eq:detformulas} to $10^{-112}$; and the index
law of Theorem~\ref{thm:cfindex} exactly on twenty synthetic single-defect
measures and four graph-Laplacian-type specimens (first anomaly
$2n_1-1$ as a length, resp.\ $2n_0-2$ as a mass, in every case).

\subsection{The a.c.\ reference at $\omega>0$: interval law at all
displacements}
For a scaled reference family (mock ordinates $\gamma_j/5$, Poisson
broadening $\omega\in\{0.15,0.3\}$) with one quartet, measured detection
orders match the interval Green prediction with no fitted constants across
effective displacements $\eta\in[0.2,1.2]$ (ratios $0.64$--$1.01$ over
eight cases), including displacements far below the ordinate spacing, where
the unbroadened (pure-point) background deviates; two parameter pairs with
equal $\eta$ and different $(\delta,\omega)$ give equal depth, confirming
that the defect enters only through the horizontal excess.

\subsection{Unscaled zeta data and the quasi-discrete regime}
With the first $100$ genuine ordinates plus the Riemann--von Mangoldt tail,
$s_0=200$, $\omega\in\{0.1,0.2\}$ (control clean to $n=28$):

\begin{center}
\begin{tabular}{cccccccc}
\toprule
$\omega$ & $\tau$ & $\eta$ & $n_0$ & $n_1$ & $g$ exact & $|\eps_{\rm eff}|$
& measured/naive \\
\midrule
0.1 & 16 & 1.0 & 14 & 14 & 0.18085 & 0.00438 & 0.93 \\
0.1 & 20 & 1.0 & 14 & 14 & 0.09379 & 0.00333 & 0.46 \\
0.1 & 25 & 1.0 & 22 & 21 & 0.05866 & 0.00242 & 0.41 \\
0.1 & 16 & 0.5 & 19 & 18 & 0.09272 & 0.00439 & 0.61 \\
0.1 & 20 & 0.5 & 26 & 25 & 0.04705 & 0.00333 & 0.41 \\
0.2 & 20 & 1.0 & 16 & 16 & 0.09379 & 0.00333 & 0.53 \\
\bottomrule
\end{tabular}
\end{center}

\noindent
Detection at genuine zeta scale is finite and computable; the sub-unity
ratios reflect the \emph{quasi-discrete regime}: at these heights the local
zero spacing $2\pi/\log(\tau/2\pi)\approx4.5$--$5.4$ exceeds both $\eta$ and
$\omega$, and the interval law overestimates the depth. In controlled
scaled experiments the ratio approaches $1$ once the displacement exceeds
roughly $1.5$--$3$ local spacings.

\subsection{$s_0$-dependence and the time-domain constant}
The $s_0$ panel ($n_{s_0}=22,14,12,17,28$ at $s_0=1,4,25,100,400$, fixed
data) documents the genuine scale dependence of
Section~\ref{sec:s0}. For the time domain, the solved envelope thresholds
converge to the constant of Theorem~\ref{thm:envelope}:
residuals $\eta t_*-2\log\tau-\log\log\tau=-3.757\to-3.483$ over
$\tau=10^2$--$10^{12}$ against $\log(S_\omega/\eta)=-3.480$ at
$\omega=0.1$, $\eta=0.15$.

%% Figures (to be generated for the submission version):
%% Fig 1: schematic — support, forbidden atom, conjugate pair, Green contours,
%%        CF chain with first bad coefficient highlighted.
%% Fig 2: eps*|a|*K'_n(a,a) vs n for a toy background (exact threshold).
%% Fig 3: table above as plot, annotated with the quasi-discrete regime.
%% Fig 4: eta*t* - 2log(tau) - loglog(tau) -> log(S_omega/eta).
%% Fig 5: n_{s0} vs s0 (non-monotone).

% ============================================================================
\section{Discussion and open problems}\label{sec:discussion}
% ============================================================================

\subsection{What has been shown}
An exact chain --- finite-rank defect $\to$ Hankel perturbation $\to$
Christoffel threshold $\to$ first continued-fraction failure --- with a
Green-function resolution law at fixed defect, and an arithmetic realization
in which the defect parameter of the shifted zeta family is the same
horizontal excess $\eta=\delta-\omega$ that drives the time-domain
criterion. The two exact representations of the same hypothetical violation
have parametrically different conditioning: logarithmic in height for the
screw function, polynomial in height (in the regimes established) for the
compactified inverse-spectral coordinate.

\subsection{What has not been shown}
No new zero-free region; no reformulation that makes the Riemann hypothesis
more accessible; no claim that an actual failure of RH would consist of a
single quartet over a positive background. The moving-point law
\eqref{eq:highheight} awaits uniform Christoffel asymptotics; the constants
in the depth law are provably not universal
(Remark~\ref{rem:noO1}, Section~\ref{sec:s0}).

\subsection{Open problems}
\begin{openproblem}[Endpoint uniformity]
As $\omega\downarrow0$ the reference measure degenerates from absolutely
continuous to pure point, and the observed behavior crosses over from the
interval law to a discrete-gap regime governed by the ratio of the defect
displacement to the local zero spacing. Prove a uniform (double-scaling)
Christoffel asymptotic covering this degeneration.
\end{openproblem}
\begin{openproblem}[Moving-point uniformity]
Establish \eqref{eq:highheight} as a theorem, e.g.\ through strong
asymptotics for the compactified $\omega$-density or extremal polynomial
bounds uniform near the endpoint.
\end{openproblem}
\begin{openproblem}[Optimal scale]
Characterize the minimizer $s_0^*(\tau)$ of $n_{s_0}$ and its asymptotics.
\end{openproblem}
Further directions: multiple-defect interaction laws
(Remark~\ref{rem:multidefect}) and their rate-tie phenomena; analogous
sensitivity laws for other completed $L$-functions, where the same
construction applies verbatim to the corresponding shifted families; and the
finite (graph zeta) laboratory in which the present mechanism was first
isolated, which raises separate questions about non-backtracking spectra and
twisted responses and will be treated elsewhere.

\subsection*{Acknowledgements}
AI tools were used extensively in exploration, numerical verification,
auditing, and drafting; the author directed the research, reviewed the
arguments and sources, and takes responsibility for the content.
% (Placeholder wording; align with the companion papers' final formulation.)

% ============================================================================
\begin{thebibliography}{99}
% ALL ENTRIES: verify bibliographic data before submission. %% VERIFY

\bibitem{Stieltjes} T.~J.~Stieltjes,
\emph{Recherches sur les fractions continues},
Ann.\ Fac.\ Sci.\ Toulouse 8 (1894), J1--J122. %% VERIFY

\bibitem{Wall} H.~S.~Wall,
\emph{Analytic Theory of Continued Fractions},
Van Nostrand, 1948. %% VERIFY chapter for determinant formulas

\bibitem{GK} F.~R.~Gantmacher and M.~G.~Kre\u{\i}n,
\emph{Oscillation Matrices and Kernels and Small Vibrations of Mechanical
Systems}, AMS Chelsea, 2002 (orig.\ 1950). %% VERIFY mechanical CF interpretation

\bibitem{KacKrein} I.~S.~Kac and M.~G.~Kre\u{\i}n,
\emph{On the spectral functions of the string},
Amer.\ Math.\ Soc.\ Transl.\ (2) 103 (1974), 19--102. %% VERIFY length=q(0) identity

\bibitem{Schmudgen} K.~Schm\"udgen,
\emph{The Moment Problem}, Springer GTM 277, 2017. %% VERIFY

\bibitem{Uvarov} V.~B.~Uvarov,
\emph{The connection between systems of polynomials orthogonal with respect
to different distribution functions},
USSR Comput.\ Math.\ Math.\ Phys.\ 9 (1969), 25--36. %% VERIFY

\bibitem{Zhedanov} A.~Zhedanov,
\emph{Rational spectral transformations and orthogonal polynomials},
J.\ Comput.\ Appl.\ Math.\ 85 (1997), 67--86. %% VERIFY

\bibitem{Nevai} P.~Nevai,
\emph{G\'eza Freud, orthogonal polynomials and Christoffel functions. A case
study}, J.\ Approx.\ Theory 48 (1986), 3--167. %% VERIFY

\bibitem{Totik} V.~Totik,
\emph{Asymptotics for Christoffel functions for general measures on the real
line}, J.\ Anal.\ Math.\ 81 (2000), 283--303. %% VERIFY

\bibitem{StahlTotik} H.~Stahl and V.~Totik,
\emph{General Orthogonal Polynomials},
Cambridge Univ.\ Press, 1992. %% VERIFY regularity criterion + root asymptotics

\bibitem{BergDurea} C.~Berg and A.~J.~Dur\'an,
\emph{A transformation from Hausdorff to Stieltjes moment sequences},
Ark.\ Mat.\ 42 (2004), 239--257. %% VERIFY

\bibitem{LasserrePauwels} J.-B.~Lasserre and E.~Pauwels,
\emph{The empirical Christoffel function with applications in data
analysis}, Adv.\ Comput.\ Math.\ 45 (2019), 1439--1468. %% VERIFY

\bibitem{ChristoffelAnomaly2026}
\emph{Scalable anomaly detection via a univariate Christoffel function},
preprint, arXiv:2606.12483 (2026). %% VERIFY authors

\bibitem{DataStreams}
\emph{Leveraging the Christoffel function for outlier detection in data
streams}, Int.\ J.\ Data Sci.\ Anal.\ (2024). %% VERIFY authors/volume

\bibitem{Suzuki2023} M.~Suzuki,
\emph{Aspects of the screw function corresponding to the Riemann
zeta-function}, J.\ London Math.\ Soc.\ 108 (2023). %% VERIFY pages

\bibitem{MarginPaper} [Author],
\emph{The linear margin of the shifted screw-function family} (companion
paper), 2026. %% placeholder — final title/reference

\bibitem{StringPaper} [Author],
\emph{The $\omega$-string family of the Riemann xi function} (companion
paper), 2026. %% placeholder — final title/reference

\bibitem{ROUND6} [Author],
manuscript-preparation audit record, 2026. %% internal; cite or absorb

\bibitem{LabRecord} [Author],
laboratory record: finite defect-depth experiments, rounds 1--6, 2026.
%% internal record; convert relevant tables into the paper or a repository link

\end{thebibliography}

\end{document}
